% -*- coding: utf-8 -*-
% !TEX program = xelatex
\documentclass[plain,noanswer,medmath]{../../template/jnuexam}
\usepackage{../../template/kysx}

\begin{document}


\examtitle{2008}{2}


\loadproblems{1}{2008P1}%载入试卷一


\makepart{选择题}{1～8小题，每小题4分，共32分}


\begin{problem}
设$f(x) = x^2(x - 1)(x - 2)$，求$f'(x)$的零点个数为\pickout{}
\begin{abcd}
\item $0$．
\item $1$．
\item $2$．
\item $3$．
\end{abcd}
\end{problem}


\begin{problem}
如图，曲线段方程为$y = f(x)$，函数在区间$[0,a]$上有连续导数，
则定积分$\int_0^a xf'(x) \dx$等于\pickout{}
\par\centerline{\includegraphics[width=0.5\linewidth]{2008B1}}
\begin{abcd}
\item 曲边梯形$ABOD$面积．
\item 梯形$ABOD$面积．
\item 曲边三角形$ACD$面积．
\item 三角形$ACD$面积．
\end{abcd}
\end{problem}


\useproblem{1}{1}{3}%【同试卷一第一(3)题】


\begin{problem}
设函数$f(x) = \frac{\ln |x|}{|x - 1|}\sin x$，则$f(x)$有\pickout{}
\begin{abcd}
\item $1$个可去间断点，$1$个跳跃间断点．
\item $1$个可去间断点，$1$个无穷间断点．
\item $2$个无穷间断点．
\item $2$个跳跃间断点．
\end{abcd}
\end{problem}


\useproblem{1}{1}{4}%【同试卷一第一(4)题】


\begin{problem}
设函数$f$连续．
若$F(u,v) = \iint_{D_{uv}} \frac{f\left(x^2 + y^2\right)}{\sqrt{x^2 + y^2}} \dx\dy$，
其中区域$D_{uv}$为图中阴影部分，则$\frac{\pdF}{\pdu} =$\pickout{}
\par\centerline{\includegraphics[width=0.5\linewidth]{2008B2}}
\begin{abcd}
\item $vf\left(u^2 \right)$	．
\item $\frac{v}{u}f\left(u^2 \right)$．
\item $vf(u)$．
\item $\frac{v}{u}f(u)$．
\end{abcd}
\end{problem}


\useproblem{1}{1}{5}%【同试卷一第一(5)题】


\begin{problem}
设$A = \begin{pmatrix}
  1 & 2 \\
  2 & 1
\end{pmatrix}$，则在实数域上与$A$合同的矩阵为\pickout{}
\begin{abcd}
\item $\begin{pmatrix}
  -2 & 1 \\
   1 & -2
\end{pmatrix}$．
\item $\begin{pmatrix}
   2 & -1 \\
  -1 & 2
\end{pmatrix}$．
\item $\begin{pmatrix}
  2 & 1 \\
  1 & 2
\end{pmatrix}$．
\item $\begin{pmatrix}
   1 & -2 \\
  -2 & 1
\end{pmatrix}$．
\end{abcd}
\end{problem}


\makepart{填空题}{9～14小题，每小题4分，共24分}


\begin{problem}
已知$f(x)$连续，且$\lim_{x \to 0} \frac{1 - \cos[xf(x)]}{(\e^{x^2} - 1)f(x)} = 1$，
则$f(0) =$ \fillin{}．
\end{problem}


\begin{problem}
微分方程$(y + x^2\e^{- x})\dx - x\dy = 0$的通解是$y =$ \fillin{}．
\end{problem}


\useproblem{1}{2}{10}%【同试卷一第二(10)题】


\begin{problem}
求函数$f(x) = (x - 5)x^{2/3}$的拐点坐标为 \fillin{}．
\end{problem}


\begin{problem}
已知$z = \left(\frac{y}{x} \right)^{\frac{x}{y}}$，
则$\left. \frac{\pdz}{\pdx} \right|_{(1,2)} =$ \fillin{}．
\end{problem}


\begin{problem}
设$3$阶矩阵$A$的特征值是$2,3,\lambda$，若行列式$|2A| = - 48$，则$\lambda =$ \fillin{}．
\end{problem}


\makepart{解答题}{15～23小题，共94分}


\useproblem[points=9]{1}{3}{15}%【同试卷一第三(15)题】


\begin{problem}[points=10]
设函数$y = y(x)$由参数方程$\begin{cases}
  x = x(t), \\
  y = \int_0^{t^2} \ln(1 + u)\du
\end{cases}$确定，其中$x(t)$是初值问题$\begin{cases}
  \frac{\dx}{\dt} - 2t\e^{- x} = 0, \\
  x|_{t = 0} = 0
\end{cases}$的解，求$\frac{\d^2y}{\dx^2}$．
\end{problem}


\begin{problem}[points=9]
计算$\int_0^1 \frac{x^2\arcsin x}{\sqrt{1 - x^2}} \dx$．
\end{problem}


\begin{problem}[points=11]
计算$\iint_D \max \left\{xy,1 \right\} \dx\dy$，
其中$D = \{\left. (x,y) \right|0 \le x \le 2,0 \le y \le 2\}$．
\end{problem}


\begin{problem}[points=11]
设$f(x)$是区间$[0, +\infty)$上具有连续导数的单调增加函数，
且$f(0) = 1$．对于任意的$t \in[0, +\infty)$，直线$x = 0$，$x = t$，
曲线$y = f(x)$以及$x$轴所围成曲边梯形绕$x$轴旋转一周生成一旋转体．
若该旋转体的侧面面积在数值上等于其体积的$2$倍，求函数$f(x)$的表达式．
\end{problem}


\begin{problem}[points=11]
\begin{enumerate}
\item 证明积分中值定理：若函数$f(x)$在闭区间$[a,b]$上连续，则至少存在一点$\eta\in[a,b]$，
      使得$\int_a^b f(x)\dx = f(\eta)(b - a)$．
\item 若函数$\varphi (x)$具有二阶导数，
      且满足$\varphi(2) > \varphi(1)$，$\varphi(2) > \int_2^3\varphi(x)\dx$，
      则至少存在一点$\xi \in(1,3)$，$\text{使得}\varphi''(\xi) < 0$．
\end{enumerate}
\end{problem}


\begin{problem}[points=11]
求函数$u = x^2 + y^2 + z^2$在约束条件$z = x^2 + y^2$和$x + y + z = 4$下的最大和最\goodbreak 小值．
\end{problem}


\useproblem[points=12]{1}{3}{21}%【同试卷一第三(21)题】


\begin{problem}[points=10]
设$A$为$3$阶矩阵，$\alpha_1$, $\alpha_2$为$A$的分别属于特征值$- 1$, $1$的特征向量，
向量$\alpha_3$满足$A\alpha_3 = \alpha_2 + \alpha_3$．\par
\begin{enumerate*}
\item 证明$\alpha_1,\alpha_2,\alpha_3$线性无关；
\item 令$P = \left(\alpha_1,\alpha_2,\alpha_3 \right)$，求$P^{-1} AP$．
\end{enumerate*}
\end{problem}


\end{document}
